\section{Methodology}
\label{sec:methodology}

\subsection{Prompt-Design Space}
\label{sec:methodology:prompts}

We study five independently toggleable prompt components layered on a
common base instruction (classify the input text as containing sexism or
not): a system-level \textbf{role} specification identifying the model as
a sexism-annotation assistant; a bulleted list of \textbf{aspects}
(seven linguistic markers associated with sexist language: objectification,
gender stereotyping, hostile language, benevolent sexism, dehumanization,
double standards, and dismissiveness); a short \textbf{context} block
naming the platform, content type, and language; \textbf{chain-of-thought}
(T) instructions asking the model to reason before answering; and
\textbf{few-shot} (F) demonstrations.

Rather than the full $2^5{=}32$-cell factorial, we evaluate a curated grid
of 18 variants (V1--V18): the bare baseline with every component off (V1),
each component in isolation (V2--V6), and twelve further combinations
spanning pairwise through five-way interactions, culminating in all five
components active at once (V18). Table~\ref{tab:main-grid}
(Sec.~\ref{sec:results:grid}) lists the exact component composition of
every variant alongside its results, so it is not repeated here.

Active (non-role) components are concatenated into the prompt's user turn
in a fixed order (context, chain-of-thought, aspects, few-shot) ahead
of the input text; the role component, when active, is instead injected as
a separate system-role chat message. Non-chain-of-thought variants append
an explicit trailing \texttt{ANSWER:} cue immediately before generation,
forcing an immediate two-token \texttt{YES}/\texttt{NO} decision;
chain-of-thought variants omit this cue and instead instruct the model to
reason first and conclude with its own \texttt{ANSWER: YES} /
\texttt{ANSWER: NO} line, so the forced cue never pre-empts the model's own
reasoning (Sec.~\ref{sec:experimental:generation} details the corresponding
generation-budget difference between the two). Few-shot variants sample
$k{=}2$ demonstrations per class (2 sexist, 2 not-sexist; 4 total),
seeded ($\text{seed}{=}42$) and drawn without replacement from the EDOS
training split (Sec.~\ref{sec:experimental:dataset}), formatted as
\texttt{TEXT}/\texttt{ANSWER} pairs.

\subsection{Model Roster}
\label{sec:methodology:models}

We evaluate six open-weight, instruction-tuned LLMs spanning five model
families and a 3.8--9B parameter range (Table~\ref{tab:models}). Every
model is loaded from a Hugging Face Hub revision pinned to a fixed commit
hash, so a later upstream update to a repository's default branch cannot
silently change which weights a rerun evaluates. Two models
(llama3.1, gemma2) are gated, requiring an accepted Hugging Face license;
the remaining four are openly accessible. qwen3 is the only reasoning-tuned
model in the roster: its chat template exposes a runtime
\texttt{enable\_thinking} toggle, wired to each variant's own
chain-of-thought flag, giving it genuine extended reasoning on
chain-of-thought variants rather than the brief prose reasoning the other
five models produce from instructions alone.

\begin{table}[t]
\centering
\footnotesize
\caption{Evaluated model roster. All Hugging Face Hub revisions are
pinned to a fixed commit (Sec.~\ref{sec:methodology:models}).}
\label{tab:models}
\begin{tabular}{lllc}
\toprule
Model & Hugging Face ID & Params & Reasoning \\
\midrule
mistral   & mistralai/Mistral-7B-Instruct-v0.3 & 7B   & \\
phi3      & microsoft/Phi-3-mini-4k-instruct   & 3.8B & \\
llama3.1  & meta-llama/Llama-3.1-8B-Instruct   & 8B   & \\
qwen2.5   & Qwen/Qwen2.5-7B-Instruct           & 7B   & \\
gemma2    & google/gemma-2-9b-it               & 9B   & \\
qwen3     & Qwen/Qwen3-8B                      & 8B   & \checkmark \\
\bottomrule
\end{tabular}
\end{table}

\subsection{Fine-Tuned and Classical Baselines}
\label{sec:methodology:baselines}

All three baselines below train on the EDOS training split and are scored
on the same held-out test split as the prompting grid
(Sec.~\ref{sec:experimental:dataset}), for direct comparability with the
prompting results (Sec.~\ref{sec:results:baselines}).

\textbf{QLoRA.} For each of the six prompted models we additionally
fine-tune a QLoRA adapter: 4-bit NF4-quantized frozen base weights (bf16
compute dtype) with LoRA adapters ($r{=}16$, $\alpha{=}32$, dropout
$0.05$) targeting the attention projection matrices
(\texttt{q\_proj}, \texttt{k\_proj}, \texttt{v\_proj}, \texttt{o\_proj}).
Training runs 6 epochs (learning rate $2{\times}10^{-4}$, batch size 16,
max sequence length 512) on the training split, oversampled to a balanced
50/50 class ratio before tokenization: the SFT objective (next-token
prediction on a single \texttt{" YES"}/\texttt{" NO"} completion) has no
built-in class weighting, and, trained directly on EDOS's natural
$\sim$76/24 imbalance, collapses to unconditionally predicting the
majority class. Both the fine-tuning prompt template and completion format
exactly match variant V1's own zero-component prompt, so the fine-tuned
model is compared like-for-like against the prompting-only results at
inference time: same instruction, same input format, only the weights
differ. We save a per-epoch adapter snapshot and promote the epoch with
the highest F1 on the EDOS dev split to the one evaluated on the test
split, mirroring the encoder baseline's own selection criterion below.

\textbf{Encoder.} A \texttt{roberta-base} sequence classifier
(2 output labels), fine-tuned for 3 epochs (learning rate
$2{\times}10^{-5}$, batch size 64, max sequence length 128) directly on
the training split, with the epoch achieving the best dev-split F1
promoted to the reported test-set number.

\textbf{Classical.} A non-neural baseline: TF-IDF features
(unigrams and bigrams, minimum document frequency 2, 50,000-feature cap,
sublinear term-frequency scaling) feeding a logistic regression classifier
(balanced class weighting to offset EDOS's natural imbalance, up to 1,000
solver iterations), fit on the same training split.

\subsection{Statistical Testing}
\label{sec:methodology:stats}

For each model we test every non-baseline variant (V2--V18) against V1
with two independent, per-model-family-corrected tests: an exact
McNemar's test on the accuracy-based discordant-pair counts between the
two variants' predictions, and a paired bootstrap test on the difference
in F1 (2,000 resamples, percentile confidence interval). Both tests'
p-values are Holm-corrected within their own family of 17 comparisons per
model, rather than pooled across the two tests. McNemar and the F1-diff
bootstrap test the same baseline-vs-variant hypothesis via different
statistics, so correcting one family with the other's p-values would
misrepresent the family-wise error control each provides on its own
(Sec.~\ref{sec:results:significance} discusses a case where the two tests
disagree).

\subsection{Sensitivity-Sweep Design}
\label{sec:methodology:sweeps}

Three supplementary robustness checks target each model's own
best-performing main-grid variant only, not the full grid along a second
axis, which would be prohibitively expensive: (i) an \textbf{order sweep},
rerendering that variant's active components under up to 5 distinct,
seeded random orderings of \{context, chain-of-thought, aspects,
few-shot\} (deduplicated to distinct induced orders, since an inactive
component does not affect ordering), holding the active component set
itself fixed; (ii) a \textbf{few-shot-size sweep}, rerunning the variant
at $k\in\{1,3,4\}$ demonstrations per class against the grid's own default
$k{=}2$, for variants where few-shot is active; and (iii) a
\textbf{demonstration-seed sweep}, redrawing which $k{=}2$-per-class
demonstrations are sampled under three alternate random seeds (1, 7, 123)
against the grid's own fixed seed (42), holding both the component order
and the demonstration count fixed, for variants where few-shot is active.
All three sweeps reuse the same evaluation engine and generation settings
as the main grid (Sec.~\ref{sec:experimental:generation});
Sec.~\ref{sec:results:sensitivity} and
Sec.~\ref{sec:results:seed-sensitivity} report the results.
